% ============================================================
% Section 2: Related Work
% ============================================================

\section{Related Work}
\label{sec:related}

\noindent\textbf{P\&ID Digitization.}
Early approaches decompose P\&ID digitization into sequential modules:
symbol detection using CNNs~\cite{rahul2019,mani2020,cha2019},
text recognition~\cite{baek2019craft}, and line detection via
Probabilistic Hough Transform~\cite{stuermer2023}.
These pipelines are brittle because errors in each stage propagate
forward, and the final graph is assembled without global structural context.
Jamieson~\etal~\cite{jamieson2024} survey the field and explicitly
identify data scarcity as the primary barrier to progress.
The first end-to-end graph-extraction approach is
Stürmer~\etal~\cite{sturmer2025}, who apply Relationformer to the
OPEN100 benchmark and achieve a 25\%+ improvement in edge detection over
the modular baseline.
Crucially, they require 60 annotated real P\&IDs for training.
We address the orthogonal question: what happens when none are available?

\noindent\textbf{Synthetic Data for Engineering Diagrams.}
Paliwal~\etal~\cite{paliwal2021} release Dataset-P\&ID, 500 synthetic
P\&IDs generated by randomly placing 32 symbol templates on a blank canvas.
Stürmer~\etal~\cite{sturmer2025} use a similar approach for pre-training
(2,000 template-synth images) but find that synth-only training
``suffers significantly'' on real diagrams (Appendix~A5).
Nurminen~\etal~\cite{nurminen2020} and Elyan~\etal~\cite{elyan2020} use
GAN-based augmentation for visual diversity, but GANs can only modify
image appearance — they cannot generate structurally new graphs.
We identify the root cause of the template failure: randomly placed symbols
produce graphs with unrealistic degree distributions inconsistent with
real plant topology.
Our topology-preserving generator addresses this directly.

\noindent\textbf{Simulation-to-Real Transfer.}
Domain adaptation from synthetic to real data is well-studied in
natural image settings~\cite{tobin2017,ganin2016}, but has received
little attention for structured diagram understanding.
The key insight from our work echoes findings in robotics sim-to-real
literature~\cite{akkaya2019}: structural fidelity of the simulated
environment matters more than visual realism for transfer.
For P\&IDs, \emph{graph topology} plays the role of physical dynamics —
a model trained on structurally unrealistic graphs fails to transfer
even when visual appearance is plausible.

\noindent\textbf{Image-to-Graph Transformers.}
Relationformer~\cite{shit2022} extends Deformable DETR~\cite{zhu2021}
with a shared relation token for joint object detection and edge
prediction.
EGTR~\cite{im2024} extracts scene graphs from DETR attention matrices
without bespoke relation tokens, achieving higher accuracy on Visual
Genome.
We adopt Relationformer following~\cite{sturmer2025} to ensure that
performance differences between our configurations are attributable to
training data rather than architecture.
